\chapter{Concpetion Hardware}

\section{Analyse fonctionnelle hardware}

\fig[H, width=1\textwidth]{Analyse fonctionnelle - Niveau 1}{Analyse_Fonctionnelle-Niv1.drawio}

\fig[H, width=0.75\textwidth]{Analyse fonctionnelle - Niveau 2}{Analyse_Fonctionnelle-Niv2.drawio}

\clearpage
\section{Choix des composants}
Cette section traitera le choix des composants. En premier lieu, le micro-contrôleur, qui est le coeur du système, sera choisi.

Ensuite, les composants externes, qui seront visibles depuis l'extérieur du boîtier et avec lesquels l'utilisateur interragira, seront sélectionnés.
Ces composants doivent être choisi en priorité pour pouvoir par la suite designer le boîtier et connaître la taille finale du PCB à réaliser.

Et finalement le restant des composants sera choisi.

\subsection{Micro-contrôleur}

\subsection{Écran}

\subsection{Batterie}
Pour déterminer la capacité de la batterie, il faut se baser sur la consommation maximale des gros consommateurs et de de l'autonomie souhaitée. Dans ce cas, les gros consommateurs sont l'écran, le micro-contrôleur, le DAC et l'amplificateur.

Voici la consommation maximale de ces composants selon le datasheet :

\begin{itemize}
    \item Écran OLED : $\sim$ 60 mA max.
    \item DAC : $\sim$ 20 mA max.
    \item Amplificateur audio : $\sim$ 25 mA max.
    \item Microcontrôleur ESP32 : $\sim$ Bonjour.
\end{itemize}

\subsection{DAC}
Le choix d'un DAC se fait principalement sur les caractéristiques suivantes :

\begin{itemize}
    \item Résolution en bit
    \item Frequence d'échantillonnage en Hz
\end{itemize}

\subsection{Boutons}

\subsection{Contrôle du volume}

\clearpage
\section{Design du boîtier}